\chapter{Literature Review}
\label{ch:literature-review}
This chapter summarizes the methodological and technical background that later chapters build on: the CRISP-DM process model that structures this entire report, and the four technical areas — PDF-based information extraction, knowledge graphs, lexical versus semantic text retrieval, and retrieval-augmented generation with large language models — that the system's design draws on.

\section{CRISP-DM as a process model}
The Cross-Industry Standard Process for Data Mining (CRISP-DM) was published in 2000 by a consortium of companies (NCR, SPSS, Daimler-Benz, and OHRA) as a vendor-neutral methodology for data mining and, more broadly, data-driven system development. It defines six phases — Business Understanding, Data Understanding, Data Preparation, Modeling, Evaluation, and Deployment — explicitly organized as an iterative cycle rather than a strict waterfall: the standard's own process diagram shows arrows looping back from later phases to earlier ones, reflecting the reality that data understanding routinely reveals gaps in the original business framing, and modeling routinely reveals gaps in data preparation that send a project back a step.

CRISP-DM was designed with classical predictive modeling in mind (a target variable, a training set, a model that scores held-out data), and this project's system is not a predictive model in that sense — it is a hybrid of a deterministic knowledge base, a lexical search index, and a constrained generative-language layer. Applying CRISP-DM's phase structure to a system of this shape required interpretation, documented explicitly at each phase boundary in this report, but the underlying discipline transfers directly and is arguably more important for a system like this one than for a classical model: understanding the data's real structure and failure modes (Data Understanding) before building anything on top of it, and evaluating the system's actual correctness — not just whether it runs — before calling it done (Evaluation), are exactly the two disciplines a project handling compliance-sensitive facts cannot afford to skip.

\section{Knowledge graphs as an enterprise data structure}
A knowledge graph represents entities as nodes and relationships between them as typed, directed edges, typically stored and queried as a graph rather than as normalized relational tables. The approach has a long history in enterprise information systems, gaining particular prominence for problems where the relationships between entities are as important as the entities themselves — exactly the shape of the DAM's own underlying structure, where the meaningful unit of information is not "this role exists" but "this role is responsible for this activity, with this specific authority code and this specific footnoted qualification." Representing the DAM as a property graph (\texttt{networkx}'s \texttt{MultiDiGraph} implementation, in this project's case) makes relationship-shaped questions — "who is connected to this activity by an approval edge" — a direct graph traversal rather than a join across several relational tables or a re-parse of a flat list on every query.

\section{Lexical versus semantic text retrieval}
Text retrieval systems broadly split into lexical methods, which match documents to a query based on shared surface vocabulary (classically, TF-IDF — Term Frequency–Inverse Document Frequency — weighting combined with cosine similarity), and semantic methods, which match based on learned vector representations (embeddings) intended to capture meaning independent of exact wording. Semantic methods generally outperform lexical methods when there is enough training or fine-tuning data to produce embeddings well-suited to the target domain's vocabulary; lexical methods remain competitive, and are markedly cheaper and more transparent, when the corpus is small, static, and shares a controlled or specialized vocabulary — which is closer to this project's situation, as detailed in the resulting design decision in Section 8.3.

\section{PDF-based information extraction}
Extracting structured information from PDF documents is a well-known hard problem in document processing, precisely because PDF is a page-description format, not a structured-data format: it encodes where each character or word is drawn on a page, not what document structure that drawing represents. Tools like \texttt{pdfplumber} (built on \texttt{pdfminer.six}) recover words and characters with their pixel coordinates and font metrics, but reconstructing table structure, reading order, multi-line cell wrapping, and typographic conventions like superscript footnote markers is left entirely to the consuming application's own heuristics. This project's Data Understanding and Data Preparation phases (Chapters 6 and 7) are, in large part, an extended case study in exactly this class of problem: the DAM's rotated column headers, its superscript authority-code decorations, and its cross-row text wrapping are all instances of visual conventions that carry real semantic meaning but are not explicitly encoded as such anywhere in the underlying PDF structure.

\section{Retrieval-augmented generation and the hallucination problem}
Large language models are well documented to "hallucinate" — to produce fluent, confident, and factually wrong output, including inventing entities, relationships, or citations that do not exist in any source they were given. Retrieval-augmented generation (RAG), as originally proposed, addresses this by retrieving relevant source passages at query time and conditioning the model's generation on them, rather than relying solely on knowledge implicit in the model's training weights. In practice, RAG substantially reduces but does not eliminate hallucination: a model can still contradict, omit, or subtly rephrase a retrieved fact in ways that change its meaning, particularly under an instruction to "phrase this naturally," which is precisely the instruction a conversational interface needs to give it. For a system whose factual claims concern who holds financial approval authority, "substantially reduces" is not a sufficient guarantee; Section 8.4 describes this project's response to that gap — an explicit, automated verification step that checks the model's output against the retrieved facts token-for-token before it is ever shown to a user, rather than trusting the retrieval step alone to make hallucination unlikely.

